\section{追う量を変える}\label{節：追う量を変える}
差を足し、倍率を掛ければ、\bookterm{sequence}{数列}を初項からたどれる。
ここからは、その仕組みが見えるように、元の\bookterm{sequence}{数列}から別の量を作る。
各操作で消えるものと、残る変化を順に見る。

  \subsection{補助数列を設計する}\label{節：補助数列の設計}
    元の\bookterm{sequence}{数列}から作った,更新を調べるための別の\bookterm{sequence}{数列}を\textbf{\bookterm{auxiliary}{補助数列}}という.
    大切なのは置き換える記号を増やすことではなく,\textbf{何を消し,どの更新を残すか}である.
    例えば、残したい形を$b_{n+1}=pb_n$や$b_{n+1}-b_n=f(n)$と決める。
    その形になる条件を、元の式から求める。

    \bookstep{付加項を消すために引く}
    \sectionref*{節：漸化式を解く前に}の例$a_{n+1}=2a_n+3$で、$b_n=a_n+c$とおくと、
    \[
      b_{n+1}=2b_n+3-c
    \]
    である.等比型にしたいので$3-c=0$を要求すると,$c=3$が決まる.
    $3$を知っているから置くのではなく,残したい更新から引く量を決めたのである.
    一般に$a_{n+1}=pa_n+f(n)$なら,同じ更新をする既知の量$g(n)$,
    すなわち
    \[
      g(n+1)=pg(n)+f(n)
    \]
    を探す.$b_n=a_n-g(n)$とすれば,二つの式を引いて
    \[
      b_{n+1}=pb_n
    \]
    となる.定数項には定数,一次式の付加項には一次式というように,
    元の式の特徴からまず候補を絞り,係数を比較する.
    その次数では条件を満たせない場合もあり,そのときは候補を選び直す.
    候補の次数や係数の決め方は\sectionref*{節：階差等比複合型}で扱う.
    同じ更新をする量があっても,すぐに求められるとは限らない.

    \bookstep{倍率をそろえるために割る}
    $a_{n+1}=p_na_n+q_n$で倍率$p_n$が添字ごとに変わるなら,
    $h_{n+1}=p_nh_n$を満たす既知の非零の\bookterm{sequence}{数列}$h_n$を探す.
    $b_n=\frac{a_n}{h_n}$と置くと,
    \[
      b_{n+1}-b_n=\frac{q_n}{h_{n+1}}
    \]
    となる.元の倍率を$1$にそろえたので,残った変化を足せばよい.
    このように基準をそろえる操作を\textbf{\bookterm{normalize}{正規化}}という.
    具体例と既知解を引く別解は\sectionref*{節：変数係数の一次漸化式}で扱う.
    引く方法と割る方法は,同じ式にも使える候補である.
    条件を確かめる手間と、変形の後に残る計算が、選ぶ目安になる。

    \bookstep{複数の項をまとめて追う}
    次の項が直前の二項から決まるなら,$b_n=a_{n+1}-\alpha a_n$など,
    複数の項を合わせた量も候補になる.
    $b_{n+1}=\beta b_n$を目指して展開し,元の\bookterm{recurrence}{漸化式}と係数を合わせることで$\alpha,\beta$を探す.
    ここでも先に目標とする更新を決め,その条件から量を作っている.
    この考え方は\sectionref*{節：隣接3項間漸化式型}や\sectionref*{節：連立型}へつながる.

    作った量については,更新の式だけでなく\bookterm{initial}{初期条件}と元への戻し方も確認する.
    $b_n=a_n-g(n)$なら$a_n=b_n+g(n)$,$b_n=\frac{a_n}{h_n}$なら$a_n=h_nb_n$である.
    複数の項をまとめた場合は,その関係と必要な\bookterm{initial}{初期条件}から元の\bookterm{sequence}{数列}を求める.
    逆数・\bookterm{logarithm}{対数}など,ほかの量を作る場合にも,定義できる条件を確認する点は共通している.

  \subsection{候補の発見と証明を分ける}

  消したい項や、そろえたい倍率を見つける。
  そのための量を選び、実際に代入する。
  変形した式から\bookterm{general}{一般項}を求め、元の量へ戻す。
  発想は形を選ぶ段階にあり、計算はその形を確かめる段階にある。
  逆数・\bookterm{logarithm}{対数}・割り算を使う前には,その操作が定義できるかも確かめる.
  候補を探す計算の途中で条件が必要だと分かった場合は,そこで条件を確認してから続ける.

  最初の数項を計算する実験も,変形の候補を見つける助けになる.
  数項の一致は予想の根拠であり,すべての項についての証明にはならない.
  予想した式が\bookterm{initial}{初期条件}と元の更新式を満たせば,順に次の項が一意に決まる\bookterm{recurrence}{漸化式}では,\bookterm{induction}{数学的帰納法}で一致を示せる.
  証明の仕組みは、\sectionref*{節：数学的帰納法}で説明した。
  帰納法は,変形から得た式を確かめるときにも,実験から得た予想を証明するときにも使える.

  \bookterm{general}{一般項}を簡単な式で表すことに加えて,正値・\bookterm{limit}{極限}・整数性などを調べる目標もある.
  \chapterref{章：実際の解き方}では方法の選択を練習し,発展演習と付録では目標の広がりを見る.
